Considere el problema de Dirichlet para el operador de Laplace en el
disco
\begin{math}
	D=
	\left\{
	\left(x,y\right)\in\mathbb{R}^{2}\mid
	x^{2}+y^{2}<a^{2}
	\right\}
\end{math}
por
\begin{equation}\label{eq:dirichletproblemcartesian}
	\begin{cases}
		\difc.L.{u}{}=
		0   &
		\text{en }D. \\
		u=g &
		\text{en }\partial D.
	\end{cases}
\end{equation}

\begin{parts}
	\part[3]

	Reescriba el operador
	\begin{math}\displaystyle
		\difc.L.{u}{}\coloneqq
		\diffp[2]{u}{x}+
		\diffp[2]{u}{y}=
		\diffp[2]{u}{r}+
		\frac{1}{r}
		\diffp{u}{r}+
		\frac{1}{r^{2}}
		\diffp[2]{u}{\theta}
	\end{math}
	en coordenadas polares $\left(r,\theta\right)$.

	\part[3]

	Reformule el problema~\eqref{eq:dirichletproblemcartesian} en
	coordenadas polares, teniendo en cuenta que
	\begin{math}
		\partial D=
		\left\{r=a\right\}\times
		\left[0,2\pi\right)
	\end{math}.
	Luego, resuelva~\eqref{eq:dirichletproblempolar} mediante el método
	de separación de variables.
	\begin{equation}\label{eq:dirichletproblempolar}
		\begin{cases}
			\difc.L.{u}{}=
			0   &
			\text{en }\left(0,a\right)\times\left[0,2\pi\right). \\
			u=g &
			\text{en }\left\{r=a\right\}\times\left[0,2\pi\right).
		\end{cases}
	\end{equation}
\end{parts}

\begin{solutionorgrid}
	\begin{parts}
		\part

		\part

		Buscamos soluciones separadas de la forma
		\begin{math}
			u\left(r,\theta\right)=
			R\left(r\right)
			\Theta\left(\theta\right)
		\end{math}.
		\begin{equation*}
			R^{\prime\prime}
			\Theta+
			\frac{1}{r}
			R^{\prime}
			\Theta+
			\frac{1}{r^{2}}
			R
			\Theta^{\prime\prime}=
			0\iff
			r^{2}\frac{R^{\prime\prime}}{R}+
			r\frac{R^{\prime}}{R}=
			-\frac{\Theta^{\prime\prime}}{\Theta}=\lambda.
		\end{equation*}

		\begin{equation*}
			\boxed{
			u\left(r,\theta\right)=
			\frac{a^{2}-r^{2}}{2\pi}
			\int_{0}^{2\pi}
			\frac{g\left(\phi\right)}{a^{2}-2ar\cos\left(\theta-\phi\right)+r^{2}}\dl\phi.
			}
		\end{equation*}
	\end{parts}
\end{solutionorgrid}

\begin{parts}
	\part[2]

	Demuestre que las siguientes desigualdades
	$\forall x\in\mathbb{R}^{d}$:
	\begin{math}
		{\left\|x\right\|}_{1}\leq
		\sqrt{d}{\left\|x\right\|}_{2}
	\end{math}
	y
	\begin{math}
		{\left\|x\right\|}_{2}\leq
		\sqrt{d}{\left\|x\right\|}_{\infty}
	\end{math}.
	% Sugerencia: Use la desigualdad de Cauchy-Schwarz.

	\part[2]

	Sea el sistema $Ax=b$, donde $A\in\mathbb{R}^{d\times d}$ es
	invertible y $b\in\mathbb{R}^{d}$.
	Sean $x^{\ast}$ la solución exacta y $\widetilde{x}$ una
	aproximación de $x^{\ast}$.
	Demuestre que
	\begin{math}\displaystyle
		\frac{\left\|b-A\widetilde{x}\right\|}{\left\|A\right\|}\leq
		\left\|x^{\ast}-\widetilde{x}\right\|\leq
		\left\|A^{-1}\right\|
		\left\|b-A\widetilde{x}\right\|
	\end{math}.

	\part[2]

	Demuestre que
	\begin{math}\displaystyle
		\frac{{\left\|x^{\ast}-\widetilde{x}\right\|}_{2}}{{\left\|x^{\ast}\right\|}_{2}}\leq
		\kappa_{2}\left(A\right)
		\frac{{\left\|b-A\widetilde{x}\right\|}_{2}}{{\left\|b\right\|}_{2}}
	\end{math},
	donde
	\begin{math}
		\kappa_{2}
		\left(A\right)\coloneqq
		{\left\|A\right\|}_{2}
			{\left\|A^{-1}\right\|}_{2}
	\end{math}.
\end{parts}

\begin{solutionorbox}
	\begin{parts}
		\part

		Sea $x=\left(x_{1},\dotsc,x_{d}\right)^{T}\in\mathbb{R}^{d}$.
		Aplique la desigualdad de Cauchy-Schwarz a los vectores
		$a=\left(\left|x_{1}\right|,\dotsc,\left|x_{d}\right|\right)^{T}$
		y $b=\left(1,\dotsc,1\right)^{T}$.
		\begin{align*}
			{\left\|x\right\|}_{1}   & =
			\left|
			\sum_{k=1}^{d}
			\left(
			\left|x_{k}\right|\cdot 1
			\right)
			\right|=
			\left|
			\left\langle
			a,b
			\right\rangle
			\right|\leq
			{\left\|a\right\|}_{2}
			{\left\|b\right\|}_{2}=
			\sqrt{\sum_{k=1}^{d}{\left|x_{k}\right|}^{2}}
			\sqrt{\sum_{k=1}^{d}1^{2}}=
			{\left\|x\right\|}_{2}\sqrt{d}. \\
			\left\|x\right\|^{2}_{2} & =
			\sum_{k=1}^{d}
				{\left|x_{k}\right|}^{2}\leq
			\sum_{k=1}^{d}
			{\left\|x\right\|}^{2}_{\infty}=
			{\left\|x\right\|}^{2}_{\infty}
			\sum_{k=1}^{d}
			1=
			{\left\|x\right\|}^{2}_{\infty}
			d.\implies
			{\left\|x\right\|}_{2}\leq
			{\left\|x\right\|}_{\infty}
			\sqrt{d}.
		\end{align*}

		\part

		Defina $e\coloneqq x^{\ast}-\widetilde{x}$ y
		$r\coloneqq b-A\widetilde{x}=Ae$.
		\begin{align*}
			\left\|r\right\| & =
			\left\|Ae\right\|\leq
			\left\|A\right\|
			\left\|e\right\|\implies
			\left\|e\right\|\geq
			\frac{\left\|r\right\|}{\left\|A\right\|}.
			\shortintertext{Pero, $e=A^{-1}r$.}
			\left\|e\right\| & =
			\left\|A^{-1}r\right\|\leq
			\left\|A^{-1}\right\|
			\left\|r\right\|.    \\
			\frac{\left\|r\right\|}{\left\|A\right\|}\leq
			\left\|e\right\|\leq
			\left\|A^{-1}\right\|
			\left\|r\right\|.
		\end{align*}
	\end{parts}
\end{solutionorbox}

Considere el problema de valor inicial / frontera que describe una
onda viajera unidireccional con rapidez
$c\in\mathbb{R}\setminus\left\{0\right\}$
\begin{equation}\label{eq:advectionPDE}
	\begin{cases}\displaystyle
		\diffp{u}{t}+
		c\diffp{u}{x}=
		0                 &
		\text{en }\left(a,b\right)\times\left(0,1\right).   \\
		u=
		u_{0}             &
		\text{en }\left(a,b\right)\times\left\{t=0\right\}. \\
		u\left(x,t\right)=
		u\left(x,t\right) &
		\text{en }\left\{a,b\right\}\times\left(0,1\right).
	\end{cases}
\end{equation}

\begin{parts}
	\part[2]

	Aproxime en~\eqref{eq:advectionPDE} la derivada espacial de $u$ por
	una diferencia finita centrada de segundo orden,
	\begin{equation*}
		\diffp.|.{u}{x}[j]^{\mkern-10mu n}=
		\frac{
			u^{n}_{j+1}-
			u^{n}_{j-1}
		}{2\Delta x}+
		\mathcal{O}{{\left(\Delta x\right)}^{2}}.
	\end{equation*}
	Discretice el espacio utilizando una malla
	\begin{math}
		\Omega_{\Delta x}\coloneqq
		\left\{
		j\Delta x\in\left[a,b\right]\mid
		j=1,\dotsc,m
		\right\}
	\end{math}
	y considerando condiciones de frontera periódicas, de modo que
	\begin{math}
		u_{0}\left(t\right)\coloneqq
		u_{m}\left(t\right)
	\end{math}
	y
	\begin{math}
		u_{m+1}\left(t\right)\coloneqq
		u_{1}\left(t\right)
	\end{math}.
	Denotando
	\begin{math}
		u_{j}\left(t\right)\coloneqq
		u\big(x_{j},t\big)
	\end{math},
	obtenga el siguiente sistema EDO en tiempo
	\begin{equation}\label{eq:odesystem}
		\diff*{
			\begin{bmatrix}
				u_{1}\left(t\right)   \\
				u_{2}\left(t\right)   \\
				\vdots                \\
				u_{m-1}\left(t\right) \\
				u_{m}\left(t\right)
			\end{bmatrix}
		}{t}=
		\dfrac{c}{2\Delta x}
		\begin{bmatrix}
			0      & -1     & 0      & \cdots & 1      \\
			1      & 0      & -1     & \ddots & \vdots \\
			0      & \ddots & \ddots & \ddots & 0      \\
			\vdots & \ddots & \ddots & 0      & -1     \\
			-1     & 0      & \cdots & 1      & 0
		\end{bmatrix}
		\begin{bmatrix}
			u_{1}\left(t\right)   \\
			u_{2}\left(t\right)   \\
			\vdots                \\
			u_{m-1}\left(t\right) \\
			u_{m}\left(t\right)
		\end{bmatrix}.
	\end{equation}

	\part[3]

	Demuestre que $\forall j=1,\dotsc,m$:
	$\operatorname{Im}\left(\lambda_{j}\right)=0$, donde $\lambda_{j}$
	es el $j$-ésimo autovalor de la matriz antisimétrica
	circulante del sistema EDO.
\end{parts}

\begin{solutionorbox}
	.
\end{solutionorbox}
